\chapter{Foundational Models on Chest X-Rays}
\label{chap:baselines}


\section{Introduction}
Its very well known that unsupervised models helps learn better image feature than supervised models. 
For CXR a vast amount of work is done in self-supervised learning like DINO and vision-language pretraining like CLIP.  
Features or embeddings extracted from these models can be used for various tasks like classification, generation etc.
Of the various open source models available in the literature we used DINO\footnote{\url{https://huggingface.co/microsoft/rad-dino}} 
 and Biovil-T\footnote{\url{https://huggingface.co/microsoft/biovil-t}} in this chapter, due to their better performance.


The above experiments are done in order to verify the performance of the features extracted from these models, to see how well the 
embeddings can be classified. These embeddings are used in conditioning diffusion models for CXR image generation.


% \section{Self-Supervised Representation Learning with DINO}
\label{sec:dino_theory}

DINO (Self-DIstillation with NO labels) is a self-supervised learning framework that learns semantic visual representations without requiring manual annotations. It is based on knowledge distillation between two networks (student and teacher) operating on different augmented views of the same image.

\subsection{Student–Teacher Framework}

Let $\mathbf{x}$ denote an input image. We generate multiple augmented views:

\begin{equation}
\mathbf{x}_1 = t_1(\mathbf{x}), \quad
\mathbf{x}_2 = t_2(\mathbf{x})
\end{equation}

where $t_1, t_2$ are stochastic augmentation functions (cropping, color jitter, etc.).

Two networks are used:

\begin{itemize}
    \item Student network $f_{\theta_s}$
    \item Teacher network $f_{\theta_t}$
\end{itemize}

Both produce feature representations:

\begin{equation}
\mathbf{z}_s = f_{\theta_s}(\mathbf{x}_1), \quad
\mathbf{z}_t = f_{\theta_t}(\mathbf{x}_2)
\end{equation}

These representations are passed through projection heads $g_s, g_t$ to obtain logits:

\begin{equation}
\mathbf{p}_s = g_s(\mathbf{z}_s), \quad
\mathbf{p}_t = g_t(\mathbf{z}_t)
\end{equation}

\subsection{Softmax with Temperature Scaling}

The outputs are converted into probability distributions:

\begin{align}
P_s(i) &= \frac{\exp(p_s^i / \tau_s)}
{\sum_j \exp(p_s^j / \tau_s)} \\
P_t(i) &= \frac{\exp((p_t^i - c_i) / \tau_t)}
{\sum_j \exp((p_t^j - c_j) / \tau_t)}
\end{align}

where:
\begin{itemize}
    \item $\tau_s, \tau_t$ are temperature parameters,
    \item $c$ is a centering vector to stabilize training.
\end{itemize}

\subsection{Cross-Entropy Self-Distillation Loss}

The student is trained to match the teacher distribution:

\begin{equation}
\mathcal{L}_{\text{DINO}} =
- \sum_i P_t(i) \log P_s(i)
\end{equation}

Unlike contrastive methods, DINO does not use negative samples. The loss encourages invariance across augmented views.

\subsection{Teacher Momentum Update}

The teacher parameters are updated using exponential moving average (EMA):

\begin{equation}
\theta_t \leftarrow m \theta_t + (1 - m)\theta_s
\end{equation}

where $m \in [0,1]$ is a momentum coefficient close to 1.

\subsection{Avoiding Collapse}

DINO avoids trivial constant solutions via:

\begin{enumerate}
    \item Temperature sharpening of teacher outputs
    \item Centering of teacher logits
    \item EMA update of teacher network
\end{enumerate}

These mechanisms ensure that learned features capture meaningful semantic structure.

\subsection{Emergent Semantic Structure}

Empirically, DINO representations exhibit:

\begin{itemize}
    \item Strong clustering properties
    \item Attention maps aligned with object boundaries
    \item Linear separability for downstream tasks
\end{itemize}

This makes DINO particularly suitable for medical image representation learning where labels may be noisy or scarce.

% \section{Vision–Language Contrastive Learning (CLIP)}
\label{sec:clip_theory}

Contrastive Language–Image Pretraining (CLIP) learns aligned representations between images and text using large-scale image–caption pairs.

\subsection{Dual Encoder Architecture}

Given an image $\mathbf{x}$ and a text description $\mathbf{t}$:

\begin{align}
\mathbf{v} &= f_\theta(\mathbf{x}) \in \mathbb{R}^d \\
\mathbf{u} &= g_\phi(\mathbf{t}) \in \mathbb{R}^d
\end{align}

where:
\begin{itemize}
    \item $f_\theta$ is an image encoder (e.g., ViT or CNN),
    \item $g_\phi$ is a text encoder (e.g., Transformer),
    \item Both outputs are normalized:
\end{itemize}

\begin{equation}
\hat{\mathbf{v}} = \frac{\mathbf{v}}{\|\mathbf{v}\|}, \quad
\hat{\mathbf{u}} = \frac{\mathbf{u}}{\|\mathbf{u}\|}
\end{equation}

\subsection{Similarity Function}

The similarity between image $i$ and text $j$ is computed via cosine similarity:

\begin{equation}
s_{ij} = \hat{\mathbf{v}}_i^\top \hat{\mathbf{u}}_j
\end{equation}

\subsection{InfoNCE Contrastive Loss}

For a batch of $N$ aligned image–text pairs, CLIP optimizes a symmetric contrastive objective:

\begin{align}
\mathcal{L}_{\text{image}} &=
- \frac{1}{N} \sum_{i=1}^N
\log
\frac{\exp(s_{ii}/\tau)}
{\sum_{j=1}^N \exp(s_{ij}/\tau)}
\\
\mathcal{L}_{\text{text}} &=
- \frac{1}{N} \sum_{i=1}^N
\log
\frac{\exp(s_{ii}/\tau)}
{\sum_{j=1}^N \exp(s_{ji}/\tau)}
\end{align}

Total loss:

\begin{equation}
\mathcal{L}_{\text{CLIP}} =
\frac{1}{2}
\left(
\mathcal{L}_{\text{image}} +
\mathcal{L}_{\text{text}}
\right)
\end{equation}

where $\tau$ is a learnable temperature parameter.

\subsection{Contrastive Alignment Interpretation}

The objective:

\begin{itemize}
    \item Pulls matched image–text pairs closer
    \item Pushes mismatched pairs apart
    \item Structures embedding space semantically
\end{itemize}

\subsection{Zero-Shot Classification}

Given class descriptions $\{\mathbf{t}_k\}$, prediction is performed by:

\begin{equation}
\hat{y} =
\arg\max_k
\; \hat{\mathbf{v}}^\top \hat{\mathbf{u}}_k
\end{equation}

This enables zero-shot classification without fine-tuning.

\subsection{Relevance to Medical Imaging}

In medical imaging:

\begin{itemize}
    \item CLIP-like models align radiographs with report text
    \item Enable zero-shot pathology classification
    \item Provide semantically structured latent spaces
\end{itemize}

Compared to VAEs, contrastive vision–language models produce embeddings that exhibit stronger pathology-based clustering due to supervision from textual semantics.

\section{Observations}

\begin{enumerate}
  \item \textbf{Similar performance on full images.} Classification performance of RAD-DINO and BioViL-T features is almost same 
  for all the pathologies.

  \item \textbf{Joint DINO and Vision-language embeddings:}  Joint (concatenated) RAD-DINO and BioViL-T features also give same performance
  on classification as above. 

  \item \textbf{Lung segmentation does not improve performance.} With assumption that there is noise associated with the 
  outside-lung regions, same experiments were performed with lung-segmented images. While RAD-DINO features doesn't show
  any difference, BioViL-T features show a significant drop in performance for almost all the pathologies. 

\end{enumerate}

\subsection{Inference and Discussion}
\begin{itemize}

  \item \textbf{BioViL-T:} Vision--language pretraining might be losing some information like relative contrast 
  while segmenting the lung region.

  \item \textbf{DINO:} DINO might have learned more local features due to image-image pretraining. 

  \item \textbf{Joint Embeddings:} The exact same performance of joint (concatenated) RAD-DINO and BioViL-T features as above, means that 
  there is no extra information in the Vision-language features than that is present in the DINO features.

%   \item \textbf{Future work:} Dino is able to extract features from local segments while BioViL-T needs the global scale.
\end{itemize}

\subsection{Results}
% =========================================================
% Tables: Foundational models results (thr = 0.5)
% File: tables/foundational_models.tex
% =========================================================

% -------------------------
% Full CXRs — RAD-DINO
% -------------------------
\begin{table}[H]
\centering
\small
\begin{tabular}{lcccc}
\toprule
\textbf{Pathology} & \textbf{Acc} & \textbf{P} & \textbf{R} & \textbf{F1} \\
\midrule
Atelectasis               & 0.792 & 0.459 & 0.212 & 0.290 \\
Cardiomegaly              & 0.766 & 0.548 & 0.233 & 0.327 \\
Consolidation             & 0.937 & 0.000 & 0.000 & 0.000 \\
Edema                     & 0.830 & 0.580 & 0.308 & 0.402 \\
EnlargedCardiomediastinum & 0.961 & 0.000 & 0.000 & 0.000 \\
Fracture                  & 0.968 & 0.000 & 0.000 & 0.000 \\
LungLesion                & 0.961 & 0.500 & 0.010 & 0.019 \\
LungOpacity               & 0.711 & 0.587 & 0.155 & 0.245 \\
NoFinding                 & 0.826 & 0.565 & 0.389 & 0.461 \\
PleuralEffusion           & 0.807 & 0.693 & 0.634 & 0.662 \\
PleuralOther              & 0.977 & 0.000 & 0.000 & 0.000 \\
Pneumonia                 & 0.896 & 0.600 & 0.006 & 0.011 \\
Pneumothorax              & 0.967 & 0.325 & 0.174 & 0.226 \\
SupportDevices            & 0.839 & 0.726 & 0.693 & 0.709 \\
\bottomrule
\end{tabular}
\caption{Per-pathology metrics on \textbf{Full CXRs} using \textbf{RAD-DINO} features (threshold = 0.5).}
\label{tab:full_dino}
\end{table}

% -------------------------
% Full CXRs — BioViL-T
% -------------------------
\begin{table}[H]
\centering
\small
\begin{tabular}{lcccc}
\toprule
\textbf{Pathology} & \textbf{Acc} & \textbf{P} & \textbf{R} & \textbf{F1} \\
\midrule
Atelectasis               & 0.802 & 0.525 & 0.110 & 0.182 \\
Cardiomegaly              & 0.769 & 0.570 & 0.220 & 0.318 \\
Consolidation             & 0.937 & 0.000 & 0.000 & 0.000 \\
Edema                     & 0.833 & 0.645 & 0.231 & 0.341 \\
EnlargedCardiomediastinum & 0.961 & 0.000 & 0.000 & 0.000 \\
Fracture                  & 0.968 & 0.000 & 0.000 & 0.000 \\
LungLesion                & 0.961 & 0.000 & 0.000 & 0.000 \\
LungOpacity               & 0.714 & 0.604 & 0.160 & 0.252 \\
NoFinding                 & 0.820 & 0.539 & 0.372 & 0.440 \\
PleuralEffusion           & 0.808 & 0.727 & 0.571 & 0.639 \\
PleuralOther              & 0.977 & 0.000 & 0.000 & 0.000 \\
Pneumonia                 & 0.895 & 0.391 & 0.017 & 0.032 \\
Pneumothorax              & 0.968 & 0.321 & 0.118 & 0.173 \\
SupportDevices            & 0.840 & 0.751 & 0.651 & 0.697 \\
\bottomrule
\end{tabular}
\caption{Per-pathology metrics on \textbf{Full CXRs} using \textbf{BioViL-T} features (threshold = 0.5).}
\label{tab:full_biovil}
\end{table}

% -------------------------
% Full CXRs — Joint Training
% -------------------------
\begin{table}[H]
\centering
\small
\begin{tabular}{lcccc}
\toprule
\textbf{Pathology} & \textbf{Acc} & \textbf{P} & \textbf{R} & \textbf{F1} \\
\midrule
Atelectasis               & 0.802 & 0.525 & 0.121 & 0.197 \\
Cardiomegaly              & 0.775 & 0.590 & 0.250 & 0.351 \\
Consolidation             & 0.937 & 0.000 & 0.000 & 0.000 \\
Edema                     & 0.832 & 0.575 & 0.372 & 0.452 \\
EnlargedCardiomediastinum & 0.961 & 0.000 & 0.000 & 0.000 \\
Fracture                  & 0.968 & 0.000 & 0.000 & 0.000 \\
LungLesion                & 0.961 & 1.000 & 0.010 & 0.020 \\
LungOpacity               & 0.711 & 0.580 & 0.160 & 0.251 \\
NoFinding                 & 0.828 & 0.569 & 0.404 & 0.473 \\
PleuralEffusion           & 0.820 & 0.720 & 0.652 & 0.685 \\
PleuralOther              & 0.977 & 0.000 & 0.000 & 0.000 \\
Pneumonia                 & 0.895 & 0.421 & 0.015 & 0.029 \\
Pneumothorax              & 0.969 & 0.400 & 0.222 & 0.286 \\
SupportDevices            & 0.850 & 0.734 & 0.736 & 0.735 \\
\bottomrule
\end{tabular}
\caption{Per-pathology metrics on \textbf{Full CXRs} using \textbf{Joint Training} (threshold = 0.5).}
\label{tab:full_joint}
\end{table}

% -------------------------
% Lung-Segmented CXRs — BioViL-T (VAL)
% -------------------------
\begin{table}[H]
\centering
\small
\begin{tabular}{lcccc}
\toprule
\textbf{Pathology} & \textbf{Acc} & \textbf{P} & \textbf{R} & \textbf{F1} \\
\midrule
Atelectasis               & 0.823 & 0.000 & 0.000 & 0.000 \\
Cardiomegaly              & 0.822 & 1.000 & 0.002 & 0.004 \\
Consolidation             & 0.962 & 0.000 & 0.000 & 0.000 \\
Edema                     & 0.891 & 0.600 & 0.009 & 0.018 \\
EnlargedCardiomediastinum & 0.970 & 0.000 & 0.000 & 0.000 \\
Fracture                  & 0.989 & 0.000 & 0.000 & 0.000 \\
LungLesion                & 0.964 & 0.000 & 0.000 & 0.000 \\
LungOpacity               & 0.812 & 0.480 & 0.021 & 0.041 \\
NoFinding                 & 0.678 & 0.618 & 0.382 & 0.472 \\
PleuralEffusion           & 0.781 & 0.565 & 0.097 & 0.166 \\
PleuralOther              & 0.993 & 0.000 & 0.000 & 0.000 \\
Pneumonia                 & 0.935 & 0.000 & 0.000 & 0.000 \\
Pneumothorax              & 0.963 & 0.000 & 0.000 & 0.000 \\
SupportDevices            & 0.796 & 0.650 & 0.348 & 0.454 \\
\bottomrule
\end{tabular}
\caption{Per-pathology metrics on \textbf{Lung-Segmented CXRs} using \textbf{BioViL-T} features (Val split, threshold = 0.5).}
\label{tab:seg_biovil}
\end{table}

% -------------------------
% Lung-Segmented CXRs — RAD-DINO (VAL)
% -------------------------
\begin{table}[H]
\centering
\small
\begin{tabular}{lcccc}
\toprule
\textbf{Pathology} & \textbf{Acc} & \textbf{P} & \textbf{R} & \textbf{F1} \\
\midrule
Atelectasis               & 0.802 & 0.420 & 0.316 & 0.361 \\
Cardiomegaly              & 0.799 & 0.412 & 0.298 & 0.346 \\
Consolidation             & 0.953 & 0.167 & 0.062 & 0.090 \\
Edema                     & 0.877 & 0.420 & 0.337 & 0.374 \\
EnlargedCardiomediastinum & 0.965 & 0.111 & 0.022 & 0.037 \\
Fracture                  & 0.983 & 0.056 & 0.029 & 0.038 \\
LungLesion                & 0.959 & 0.200 & 0.046 & 0.075 \\
LungOpacity               & 0.753 & 0.329 & 0.307 & 0.318 \\
NoFinding                 & 0.752 & 0.661 & 0.703 & 0.681 \\
PleuralEffusion           & 0.821 & 0.609 & 0.566 & 0.587 \\
PleuralOther              & 0.993 & 0.500 & 0.045 & 0.083 \\
Pneumonia                 & 0.916 & 0.128 & 0.052 & 0.074 \\
Pneumothorax              & 0.956 & 0.339 & 0.188 & 0.241 \\
SupportDevices            & 0.835 & 0.669 & 0.636 & 0.652 \\
\bottomrule
\end{tabular}
\caption{Per-pathology metrics on \textbf{Lung-Segmented CXRs} using \textbf{RAD-DINO} features (Val split, threshold = 0.5).}
\label{tab:seg_dino}
\end{table}





